\subsection{Problem Statement}

Currently, while schools widely use digital databases to manage academic records, existing certificate verification systems still rely heavily on centralized agencies and manual authentication processes. These traditional methods present several significant challenges, including vulnerabilities to certificate forgery, inefficient verification processes, risks of manipulation of centralized data, and limited transparency in authentication procedures. Furthermore, existing systems often struggle to support the rapid and reliable verification of student credentials when studying abroad, causing delays for employers, universities, and organizations requiring verification of overseas academic qualifications. Reliance on centralized databases also increases the risk of data loss, cyberattacks, and internal manipulation. Furthermore, traditional digital certificates cannot effectively provide immutable ownership verification or ensure document integrity in a decentralized environment. As a result, organizations may struggle to establish trust in the authenticity of submitted qualifications.\\

Thus, a secure, decentralized, transparent, and tamper-proof academic certificate verification system is needed, capable of guaranteeing both document authenticity and ownership verification. Blockchain technology integrated with NFT-based digital certificates offers a potential solution to address these challenges and enhance the reliability of academic qualification verification systems.\\